\section{Linguaggi e funzioni}

\subsection{Linguaggi ricorsivi}
$L \subseteq \Sigma^*$ é \textbf{ricorsivo} sse \[\exists M :  M(x) =
\begin{cases}
	"yes" \qquad se \qquad x \in L \\
	"no" \qquad altrimenti
\end{cases}
\] \\
$M$ \textbf{decide} $L$. \\

\subsection{Linguaggi ricorsivamente enumerabili}
$L \subseteq \Sigma^*$ é \textbf{ricorsivamente enumerabile} sse \[\exists M :  M(x) =
\begin{cases}
	"yes" \qquad se \qquad x \in L \\
	\nearrow \qquad altrimenti
\end{cases}
\] \\
$M$ \textbf{accetta} $L$. \\

\subsection{Proposizione 2.1: i linguaggi ricorsivi sono ricorsivamente enumerabili}
\boxed{L \quad \textbf{ricorsivo} \implies L \quad \textbf{ricorsivamente enumerabile}} \\

\textbf{Prova:} \\
\[\exists M :  M(x) =
\begin{cases}
	"yes" \qquad se \qquad x \in L \\
	"no" \qquad altrimenti
\end{cases}
\] \\
Costruiamo $M' = (Q', \Sigma, \delta', s)$ aggiungendo un nuovo stato $q'$ che, nella nuova funzione di transizione $\delta'$, prende il posto dello stato terminale $"no"$. La funzione di transizione é uguale a quella di M in tutti gli altri casi. Il nuovo stato $q'$ sposta la testina a destra all'infinito. Formalmente:
\begin{itemize}
	\item $Q' = Q \cup \{q'\}$ \quad dove \quad $q' \notin Q$
	\item $\delta(q, \sigma) = ("no",  \tau, d) \implies \delta'(q, \sigma) = (q', \tau, d)  \qquad \forall q \in Q;  \forall \sigma, \tau \in \Gamma;  \forall d \in D$
	\item $\delta(q, \sigma) = (p,  \tau, d) \implies \delta'(q, \sigma) = (p, \tau, d)  \quad \forall q, p \in Q \setminus \{"no"\};  \forall \sigma, \tau \in \Gamma;  \forall d \in D$
	\item $\delta'(q', \sigma) = (q', \sigma, \rightarrow) \qquad \forall \sigma \in \Sigma \cup \{\sqcup\}$
\end{itemize}
\[M'(x) =
\begin{cases}
	"yes" \qquad se \qquad x \in L \\
	\nearrow \qquad altrimenti
\end{cases}
\] \\
\textbf{QED}

\subsection{Funzioni ricorsive}
La funzione totale $f : \Sigma^* \rightarrow {\Sigma_\sqcup}^*$ é \textbf{ricorsiva} sse \[\exists M :  \forall x \in \Sigma^* : M(x) = f(x) \] \\
$M$ \textbf{calcola} $f$. \\

\subsection{Funzioni ricorsivamente enumerabili}
La funzione parziale $f : \Sigma^* \rightarrow {\Sigma_\sqcup}^*$ é \textbf{ricorsivamente enumerabile} sse \[\exists M :  \forall x \in \Sigma^* : \begin{cases}
	M(x) = y \qquad se \quad f(x) = y \\
	M(x) = \nearrow \qquad altrimenti
\end{cases}
\] \\
Ricordando che le funzioni parziali non sono definite su tutto il dominio, e dunque esistono stringhe $x \in \Sigma^*$ per cui $f$ non é definita.

\subsection{Funzioni di complessitá proprie}
Le funzioni di complessitá proprie sono funzioni sui naturali, non decrescenti e tali che una macchina di Turing puó calcolarne il valore in unario in base alla lunghezza dell'input. \\
Cioé, dato un input $x$ lungo $n$, la macchina di Turing scrive tanti simboli $\sqcap$ ("quasi-blank") quanto vale $f(n)$. \\
Il tempo impiegato deve essere $O(n + f(n))$. \\
Lo spazio su ogni nastro deve essere $O(f(n))$. \\
Formalmente: una funzione $f: \mathbb{N} \rightarrow \mathbb{N}$ é \textit{propria} sse:
\begin{itemize}
	\item $f$ non é decrescente: \quad $\forall n \in \mathbb{N} \quad f(n+1) > f(n)$
	\item Esiste una MdT con input e output $M_f$ tale che, per ogni $n$: \\
		\[
		\configurazionestart{x} \xrightarrow{{M_f}^t} \configurazionepropria{x}{f(n)}
		\]
	\item $t = O(n + f(n))$
	\item $z_i = O(f(n)) \quad \forall i \in \{2, \dots, k-1\}$
\end{itemize}

Alcune funzioni proprie che useremo sono:
\begin{itemize}
	\item Costanti
	\item $\ceil{log(n)}$
	\item Polinomi
	\item Esponenziali
\end{itemize}

Le funzioni proprie sono le uniche funzioni che possono essere usate per misurare la complessitá di tempo e di spazio di un problema senza incorrere in problemi di misurazione.\\ 
Ad esempio, il Gap Theorem dimostra che esistono funzioni che "rompono" l'analisi di complessitá.

\subsubsection{Gap Theorem}
Esiste una funzione ricorsiva $f: \mathbb{N} \rightarrow \mathbb{N}$ tale che la classe di problemi risolvibili in tempo $f(n)$ coincide con quella di problemi risolvibili in tempo $2^{f(n)}$. \footnote{Vedremo meglio l'enunciato una volta introdotta la classe di complessitá \TIME{}}
